\documentclass[aspectratio=169, 14pt]{beamer}
\usepackage{stampfly_slides}

\lesson{9}{姿勢推定}{Attitude Estimation}

\begin{document}

% ------------------------------------------------------------------
\begin{frame}
  \titlepage
\end{frame}

% ------------------------------------------------------------------
\begin{frame}{今日のゴール / Today's Goal}
  \begin{block}{目標}
    相補フィルタを実装し、ESKFの推定値と比較する
  \end{block}

  \vspace{1em}

  \begin{itemize}
    \item ジャイロ積分のドリフト問題を理解
    \item 相補フィルタで加速度センサとジャイロを融合
    \item Teleplot で CF vs ESKF をリアルタイム比較
  \end{itemize}
\end{frame}

% ------------------------------------------------------------------
% Frame 3: Euler Angle Representation
% ------------------------------------------------------------------
\begin{frame}{姿勢の表現 --- ZYX オイラー角}
  \framesubtitle{Attitude Representation --- ZYX Euler Angles}

  \begin{block}{オイラー角とは}
    3次元の姿勢（回転）を 3 つの角度で表現する方法の一つ
  \end{block}

  \vspace{0.3em}

  {\small
  \begin{tabular}{lcll}
    \hline
    \textbf{記号} & \textbf{名称} & \textbf{回転軸} & \textbf{意味} \\
    \hline
    $\phi$   & Roll（ロール）  & $X''$ 軸 & 左右の傾き \\
    $\theta$ & Pitch（ピッチ） & $Y'$ 軸  & 前後の傾き \\
    $\psi$   & Yaw（ヨー）     & $Z$ 軸   & 機首方向 \\
    \hline
  \end{tabular}
  }

  \vspace{0.5em}

  \begin{alertblock}{StampFly の規約: ZYX オイラー角}
    {\small
    回転順序: $Z$ 軸 $(\psi)$ → $Y'$ 軸 $(\theta)$ → $X''$ 軸 $(\phi)$\\[2pt]
    航空宇宙分野で標準的な \textbf{Tait--Bryan 角}。ジンバルロック ($\theta = \pm 90°$) に注意。
    }
  \end{alertblock}
\end{frame}

% ------------------------------------------------------------------
% Frame 4: Euler Angle Kinematics
% ------------------------------------------------------------------
\begin{frame}{オイラー角微分とジャイロ角速度の関係}
  \framesubtitle{Euler Angle Rates vs Body Angular Velocity}

  \begin{block}{ジャイロの測定値 $\boldsymbol{\omega}$ $\neq$ オイラー角の時間微分 $\dot{\boldsymbol{\Theta}}$}
    {\small
    ジャイロはボディ座標系の角速度を測定。オイラー角の微分に変換するには変換行列 $\mathbf{E}$ が必要:
    }
  \end{block}

  {\footnotesize
  \[
  \begin{bmatrix} \dot{\phi} \\ \dot{\theta} \\ \dot{\psi} \end{bmatrix}
  =
  \underbrace{
  \begin{bmatrix}
  1 & \sin\phi\tan\theta & \cos\phi\tan\theta \\[2pt]
  0 & \cos\phi & -\sin\phi \\[2pt]
  0 & \frac{\sin\phi}{\cos\theta} & \frac{\cos\phi}{\cos\theta}
  \end{bmatrix}
  }_{\displaystyle \mathbf{E}(\phi,\,\theta)}
  \begin{bmatrix} \omega_x \\ \omega_y \\ \omega_z \end{bmatrix}
  \qquad
  \Longrightarrow
  \qquad
  \boldsymbol{\Theta}_{k+1} = \boldsymbol{\Theta}_k + \mathbf{E}\,\boldsymbol{\omega}\,\Delta t
  \]
  }

  \vspace{-0.3em}

  \begin{exampleblock}{小角度近似（$\phi,\,\theta \approx 0$）→ 相補フィルタで使用}
    {\small
    $\mathbf{E} \approx \mathbf{I}$ → $\dot{\phi} \approx \omega_x,\;\dot{\theta} \approx \omega_y,\;\dot{\psi} \approx \omega_z$
    → コード: \code{cf\_roll += gx * dt}
    }
  \end{exampleblock}
\end{frame}

% ------------------------------------------------------------------
% Frame 5: Euler Angles from Accelerometer
% ------------------------------------------------------------------
\begin{frame}{加速度からのオイラー角算出}
  \framesubtitle{Euler Angles from Accelerometer}

  \begin{block}{原理: 静止時、加速度センサは重力のみを測定}
    {\footnotesize
    ZYX 回転行列で重力ベクトルをボディ座標系に変換:
    $a_x = -g\sin\theta,\quad
    a_y = g\sin\phi\cos\theta,\quad
    a_z = g\cos\phi\cos\theta$
    }
  \end{block}

  \vspace{0.2em}

  {\small
  \begin{tabular}{lll}
    \hline
    \textbf{角度} & \textbf{計算式} & \textbf{コード} \\
    \hline
    Roll $\phi$ & $\mathrm{atan2}(a_y,\; a_z)$ & \code{atan2f(ay, az)} \\
    Pitch $\theta$ & $\mathrm{atan2}\!\bigl(-a_x,\; \sqrt{a_y^2 + a_z^2}\bigr)$ & \code{atan2f(-ax, az)}$^*$ \\
    Yaw $\psi$ & \textcolor{sfred}{算出不可}（重力は鉛直方向のみ） & 磁気センサが必要 \\
    \hline
  \end{tabular}
  }

  \vspace{0.3em}

  {\footnotesize
    \textbf{長所:} ドリフトなし（重力基準）\qquad
    \textbf{短所:} 振動・機体加速でノイズ大
  }

  \vspace{0.2em}

  {\scriptsize $^*$ 実装では小角度近似 \code{atan2f(-ax, az)} を使用（$\phi \approx 0$ のとき $\sqrt{a_y^2+a_z^2} \approx a_z$）}
\end{frame}

% ------------------------------------------------------------------
\begin{frame}{ジャイロドリフト問題 / Gyro Drift Problem}
  \begin{columns}[c]
    \begin{column}{0.48\textwidth}
      \centering
      \resizebox{\linewidth}{!}{\documentclass[border=10pt]{standalone}
\usepackage{tikz}
\usetikzlibrary{arrows.meta, decorations.pathmorphing}

\begin{document}

\definecolor{truecolor}{RGB}{0, 120, 200}     % Blue - true value
\definecolor{driftcolor}{RGB}{220, 50, 30}     % Red - gyro integrated
\definecolor{gridcolor}{RGB}{200, 200, 200}
\definecolor{labelcolor}{RGB}{80, 80, 80}
\definecolor{annocolor}{RGB}{40, 160, 40}      % Green - annotations

\begin{tikzpicture}[
    axis/.style={-{Stealth[length=3mm]}, line width=1pt, labelcolor},
    trueline/.style={truecolor, line width=2.5pt},
    driftline/.style={driftcolor, line width=2pt},
    anno/.style={annocolor, line width=1pt, dashed},
    dim/.style={{Stealth[length=2mm]}-{Stealth[length=2mm]}, annocolor, line width=0.8pt},
    annolbl/.style={font=\small\bfseries, annocolor},
  ]

  % Grid
  \foreach \x in {0, 1, ..., 12} {
    \draw[gridcolor] (\x, 0) -- (\x, 6);
  }
  \foreach \y in {0, 1, ..., 6} {
    \draw[gridcolor] (0, \y) -- (12, \y);
  }

  % Axes
  \draw[axis] (-0.3, 0) -- (12.5, 0) node[right, font=\bfseries] {$t$ [s]};
  \draw[axis] (0, -0.3) -- (0, 6.5) node[above, font=\bfseries] {$\theta$ [deg]};

  % Time labels
  \foreach \x/\t in {0/0, 2/5, 4/10, 6/15, 8/20, 10/25, 12/30} {
    \node[font=\scriptsize, below, labelcolor] at (\x, -0.1) {\t};
  }

  % True value (horizontal line at y=3)
  \draw[trueline] (0, 3) -- (12, 3);
  \node[font=\small\bfseries, truecolor, above] at (2, 3.1) {True angle};

  % Gyro integrated value (drifting upward with noise)
  \draw[driftline] plot[smooth, tension=0.5] coordinates {
    (0, 3.0) (0.5, 3.05) (1, 3.15) (1.5, 3.2) (2, 3.35)
    (2.5, 3.5) (3, 3.55) (3.5, 3.7) (4, 3.85) (4.5, 3.95)
    (5, 4.1) (5.5, 4.2) (6, 4.35) (6.5, 4.5) (7, 4.6)
    (7.5, 4.75) (8, 4.85) (8.5, 4.95) (9, 5.1) (9.5, 5.2)
    (10, 5.3) (10.5, 5.4) (11, 5.5) (11.5, 5.55) (12, 5.65)
  };
  \node[font=\small\bfseries, driftcolor, below right] at (8.5, 5.3) {Gyro $\int\omega\,dt$};

  % Drift annotation arrow at t=25s (x=10)
  \draw[anno] (10, 3.0) -- (10, 5.3);
  \draw[dim] (10.5, 3.0) -- (10.5, 5.3);
  \node[annolbl, right] at (10.6, 4.15) {Drift};

  % Small drift at t=10s (x=4)
  \draw[anno] (4, 3.0) -- (4, 3.85);
  \draw[dim] (4.5, 3.0) -- (4.5, 3.85);
  \node[annolbl, right, font=\scriptsize\bfseries] at (4.6, 3.42) {Small};

  % Bias label
  \node[font=\footnotesize, driftcolor, align=center] at (6, 0.8)
    {Gyro bias $\to$ angle error grows over time};

\end{tikzpicture}
\end{document}
}
    \end{column}
    \begin{column}{0.50\textwidth}
      \begin{itemize}
        \item ジャイロは\textbf{角速度} $\omega$ を測定
        \item 角度 $\theta = \int \omega\,dt$（積分が必要）
        \item バイアスが蓄積 → ドリフト増大
      \end{itemize}
      \vspace{0.5em}
      \begin{alertblock}{問題}
        ジャイロ単体では長時間の姿勢推定が不可能
      \end{alertblock}
    \end{column}
  \end{columns}
\end{frame}

% ------------------------------------------------------------------
\begin{frame}{相補フィルタ / Complementary Filter}
  \begin{columns}[c]
    \begin{column}{0.50\textwidth}
      \centering
      \resizebox{\linewidth}{!}{\documentclass[border=10pt]{standalone}
\usepackage{tikz}
\usetikzlibrary{arrows.meta, positioning, shapes.geometric, calc}

\begin{document}

\definecolor{sfblue}{RGB}{0, 120, 200}
\definecolor{sfred}{RGB}{220, 50, 30}
\definecolor{sfgreen}{RGB}{40, 160, 40}
\definecolor{sfgray}{RGB}{80, 80, 80}
\definecolor{sflight}{RGB}{220, 240, 255}

\begin{tikzpicture}[
    node distance=1.8cm and 2.2cm,
    block/.style={draw, sfblue, line width=1.2pt, fill=sflight,
                  minimum width=2.2cm, minimum height=1cm,
                  font=\small\bfseries, rounded corners=3pt},
    gain/.style={draw, sfblue, line width=1.2pt, fill=sflight,
                 circle, minimum size=0.9cm,
                 font=\small\bfseries},
    sum/.style={draw, sfblue, line width=1.2pt, fill=white,
                circle, minimum size=0.8cm,
                font=\small\bfseries},
    input/.style={font=\small\bfseries, sfgray},
    output/.style={font=\small\bfseries, sfblue},
    line/.style={-{Stealth[length=2.5mm]}, line width=1pt, sfgray},
    gyro/.style={sfred, line width=1.5pt},
    accel/.style={sfgreen, line width=1.5pt},
    label/.style={font=\scriptsize, sfgray},
  ]

  % --- Input nodes ---
  \node[input] (gyro_in) {Gyro $\omega$};
  \node[input, below=3.5cm of gyro_in] (accel_in) {Accel $(a_x, a_y, a_z)$};

  % --- Gyro path (top) ---
  \node[block, right=1.8cm of gyro_in, draw=sfred, fill=sfred!8] (integrator) {$\displaystyle\int dt$};
  \node[gain, right=1.8cm of integrator, draw=sfred, fill=sfred!8] (alpha) {$\alpha$};

  % --- Accel path (bottom) ---
  \node[block, right=1.8cm of accel_in, draw=sfgreen, fill=sfgreen!8] (atan) {$\mathrm{atan2}$};
  \node[gain, right=1.8cm of atan, draw=sfgreen, fill=sfgreen!8] (one_alpha) {$1\!-\!\alpha$};

  % --- Summation ---
  \node[sum, right=1.8cm of $(alpha)!0.5!(one_alpha)$] (sumnode) {$+$};

  % --- Output ---
  \node[output, right=1.5cm of sumnode] (output) {$\hat{\theta}$};

  % --- Connections ---
  % Gyro path
  \draw[line, sfred] (gyro_in) -- (integrator);
  \draw[line, sfred] (integrator) -- (alpha);
  \draw[line, sfred] (alpha) -- (alpha -| sumnode.west) -- (sumnode.west |- alpha) ;
  \draw[line, sfred] (alpha.east) -- ++(0.6,0) |- (sumnode.150);

  % Accel path
  \draw[line, sfgreen] (accel_in) -- (atan);
  \draw[line, sfgreen] (atan) -- (one_alpha);
  \draw[line, sfgreen] (one_alpha.east) -- ++(0.6,0) |- (sumnode.210);

  % Output
  \draw[line, sfblue, line width=1.5pt] (sumnode) -- (output);

  % --- Feedback loop (previous estimate fed back to integrator) ---
  \draw[line, sfgray, dashed] (sumnode.north) -- ++(0, 1.0)
    -| node[label, above, pos=0.25] {$\hat{\theta}_{k-1}$} (integrator.north);

  % --- Path labels ---
  \node[font=\scriptsize\bfseries, sfred, above=0.3cm of integrator]
    {High-pass (gyro)};
  \node[font=\scriptsize\bfseries, sfgreen, below=0.3cm of atan]
    {Low-pass (accel)};

  % --- Alpha value annotation ---
  \node[font=\footnotesize, sfgray, align=center, below=0.5cm of sumnode]
    {$\alpha = 0.98$\\[2pt]Gyro 98\% + Accel 2\%};

  % --- Output label ---
  \node[font=\footnotesize\bfseries, sfblue, above=0.15cm of output]
    {Estimated angle};

\end{tikzpicture}
\end{document}
}
    \end{column}
    \begin{column}{0.48\textwidth}
      \begin{block}{数式}
        $\hat{\theta}_k = \alpha\,(\hat{\theta}_{k-1} + \omega\,\Delta t) + (1-\alpha)\,\theta_{\mathrm{accel}}$
      \end{block}

      \vspace{0.3em}

      \begin{itemize}
        \item $\alpha = 0.98$（ジャイロ 98\% + 加速度 2\%）
        \item ジャイロ: 短期精度◎ / 長期ドリフト✕
        \item 加速度: 長期安定◎ / 振動ノイズ✕
      \end{itemize}
    \end{column}
  \end{columns}
\end{frame}

% ------------------------------------------------------------------
\begin{frame}{姿勢推定 API}
  \begin{table}
    \centering\small
    \begin{tabular}{lll}
      \hline
      \textbf{関数} & \textbf{説明} & \textbf{単位} \\
      \hline
      \api{gyro\_x/y/z()}  & 角速度（Roll/Pitch/Yaw） & rad/s \\
      \api{accel\_x/y/z()} & 加速度（X/Y/Z）          & m/s\textsuperscript{2} \\
      \api{estimated\_roll/pitch/yaw()} & ESKF 推定姿勢角 & rad \\
      \api{print(fmt, ...)} & シリアル出力（Teleplot 対応） & --- \\
      \hline
    \end{tabular}
  \end{table}

  \vspace{0.5em}

  \begin{block}{実習で使う組み合わせ}
    ジャイロ + 加速度 → 相補フィルタで姿勢を自作\\
    \api{estimated\_roll()} → ESKF の推定値と比較
  \end{block}
\end{frame}

% ------------------------------------------------------------------
\begin{frame}[fragile]{実習: 相補フィルタ + Teleplot / Hands-on}
  \begin{lstlisting}[style=sfcpp, basicstyle=\ttfamily\scriptsize, title={user\_code.cpp (excerpt)}]
#include "workshop_api.hpp"
#include <cmath>
static float cf_roll = 0.0f, cf_pitch = 0.0f;
void setup() { ws::print("Lesson 9: Attitude Estimation"); }
void loop_400Hz(float dt) {
    float gx = ws::gyro_x(), gy = ws::gyro_y();
    float ax = ws::accel_x(), ay = ws::accel_y(),
          az = ws::accel_z();
    float accel_roll  = atan2f(ay, az);
    float accel_pitch = atan2f(-ax, az);
    constexpr float alpha = 0.98f;
    cf_roll  = alpha*(cf_roll + gx*dt) + (1-alpha)*accel_roll;
    cf_pitch = alpha*(cf_pitch + gy*dt) + (1-alpha)*accel_pitch;
    // Teleplot output (VSCode Teleplot extension)
    ws::print(">cf_roll:%.1f", cf_roll*57.3f);
    ws::print(">eskf_roll:%.1f", ws::estimated_roll()*57.3f);
}
  \end{lstlisting}
\end{frame}

% ------------------------------------------------------------------
\begin{frame}[fragile]{Teleplot によるリアルタイム可視化}
  \begin{columns}[c]
    \begin{column}{0.55\textwidth}
      \begin{block}{Teleplot フォーマット}
        \api{>変数名:値} をシリアルに出力するだけ
      \end{block}
      \begin{lstlisting}[style=sfcpp, basicstyle=\ttfamily\scriptsize]
// Teleplot format: >name:value
ws::print(">cf_roll:%.2f", cf_roll * 57.3f);
ws::print(">cf_pitch:%.2f", cf_pitch * 57.3f);
ws::print(">eskf_roll:%.2f",
           ws::estimated_roll() * 57.3f);
      \end{lstlisting}
    \end{column}
    \begin{column}{0.43\textwidth}
      \begin{exampleblock}{セットアップ}
        \begin{enumerate}
          \item VSCode 拡張: \texttt{alexnesnes.teleplot} をインストール
          \item Teleplot パネルでシリアルポートを選択
          \item 自動でグラフ化される
        \end{enumerate}
      \end{exampleblock}
    \end{column}
  \end{columns}

  \vspace{0.5em}
  {\footnotesize \textbf{注意:} 400\,Hz 全 tick で出力するとシリアル帯域に負荷。4\,tick 毎（100\,Hz）にデシメーション推奨}
\end{frame}

% ------------------------------------------------------------------
\begin{frame}{Madgwick フィルタ / Madgwick Filter}
  \begin{block}{勾配降下法によるクォータニオン推定}
    チューニングパラメータが $\beta$ の \textbf{1つだけ} --- シンプルかつ高精度
  \end{block}

  \vspace{0.3em}

  \begin{columns}[c]
    \begin{column}{0.48\textwidth}
      {\small
      \begin{itemize}
        \item 状態: $\mathbf{q} = [q_0,\, q_1,\, q_2,\, q_3]$
        \item 予測: $\dot{\mathbf{q}} = \tfrac{1}{2}\,\mathbf{q} \otimes \boldsymbol{\omega}$
        \item 補正勾配:\\[2pt]
              $\nabla f = \mathbf{J}^T(\hat{\mathbf{q}},\,\hat{\mathbf{d}})\;\mathbf{f}(\hat{\mathbf{q}},\,\hat{\mathbf{d}})$
      \end{itemize}
      }
    \end{column}
    \begin{column}{0.48\textwidth}
      \begin{alertblock}{更新則}
        $\mathbf{q}_{t+1} = \mathbf{q}_t + \Delta t\Bigl(\tfrac{1}{2}\,\mathbf{q}_t \otimes \boldsymbol{\omega} - \beta\,\dfrac{\nabla f}{|\nabla f|}\Bigr)$
      \end{alertblock}
    \end{column}
  \end{columns}

  \vspace{0.3em}
  {\footnotesize $\mathbf{f}$: 加速度/磁気の目的関数、$\mathbf{J}$: そのヤコビアン、$\beta$: 補正ゲイン（$\approx 0.04$）}
\end{frame}

% ------------------------------------------------------------------
\begin{frame}{EKF（拡張カルマンフィルタ）}
  \begin{block}{非線形カルマンフィルタの5ステップ}
    状態ベクトル例: $\mathbf{x} = [\phi,\;\theta,\;\psi,\;b_{gx},\;b_{gy},\;b_{gz}]^T$
  \end{block}

  \vspace{0.3em}

  {\small
  \begin{tabular}{lll}
    \hline
    \textbf{ステップ} & \textbf{数式} & \textbf{意味} \\
    \hline
    状態予測 & $\hat{\mathbf{x}}^- = f(\hat{\mathbf{x}},\,\mathbf{u})$ & ジャイロで姿勢を積分 \\
    共分散予測 & $\mathbf{P}^- = \mathbf{F}\mathbf{P}\mathbf{F}^T + \mathbf{Q}$ & 不確かさの伝播 \\
    カルマンゲイン & $\mathbf{K} = \mathbf{P}^-\mathbf{H}^T(\mathbf{H}\mathbf{P}^-\mathbf{H}^T + \mathbf{R})^{-1}$ & 予測と観測の信頼度配分 \\
    状態更新 & $\hat{\mathbf{x}} = \hat{\mathbf{x}}^- + \mathbf{K}(\mathbf{z} - h(\hat{\mathbf{x}}^-))$ & 観測で補正 \\
    共分散更新 & $\mathbf{P} = (\mathbf{I} - \mathbf{K}\mathbf{H})\mathbf{P}^-$ & 不確かさを縮小 \\
    \hline
  \end{tabular}
  }
\end{frame}

% ------------------------------------------------------------------
\begin{frame}{ESKF（誤差状態カルマンフィルタ）}
  \begin{block}{StampFly で実際に使われている 15 状態 ESKF}
    名目状態 + 誤差状態の分離アーキテクチャ
  \end{block}

  \vspace{0.3em}

  {\small
  \begin{columns}[t]
    \begin{column}{0.48\textwidth}
      \textbf{誤差状態ベクトル（15 states）:}
      \[
        \delta\mathbf{x} = [\delta\mathbf{p},\;\delta\mathbf{v},\;\delta\boldsymbol{\theta},\;\delta\mathbf{b}_g,\;\delta\mathbf{b}_a]^T
      \]
      \begin{itemize}
        \item 予測: 名目状態はジャイロ/加速度で直接積分
        \item 誤差共分散のみ伝播
      \end{itemize}
    \end{column}
    \begin{column}{0.48\textwidth}
      \textbf{観測更新:}
      \begin{itemize}
        \item 気圧 → 高度補正
        \item ToF → 対地高度補正
        \item 磁気 → ヨー補正
        \item 光学フロー → 速度補正
      \end{itemize}
    \end{column}
  \end{columns}
  }

  \vspace{0.3em}
  {\footnotesize \textbf{利点:} 誤差状態は常に小さい → 線形化が正確}
\end{frame}

% ------------------------------------------------------------------
\begin{frame}{推定手法の比較 / Estimator Comparison}
  \centering
  \resizebox{0.95\linewidth}{!}{%
  \begin{tabular}{lcccc}
    \hline
    & \textbf{相補フィルタ} & \textbf{Madgwick} & \textbf{EKF} & \textbf{ESKF} \\
    \hline
    計算量       & 極小           & 小             & 中           & 中〜大 \\
    パラメータ数  & 1（$\alpha$）  & 1（$\beta$）   & $\mathbf{Q},\mathbf{R}$ 行列 & $\mathbf{Q},\mathbf{R}$ 行列 \\
    推定対象     & Roll/Pitch     & R/P/Yaw        & 姿勢+バイアス & 位置/速度/姿勢/バイアス \\
    使用センサ   & Gyro+Accel     & Gyro+Accel+Mag & Gyro+Accel+Mag & 全センサ（6種） \\
    精度         & ○              & ◎             & ◎            & ◎◎ \\
    \hline
  \end{tabular}
  }

  \vspace{0.8em}

  {\footnotesize
    \textbf{StampFly のデフォルト:} ESKF（15 状態）を使用。
    このレッスンでは相補フィルタを自作し、ESKF と比較して理解を深める。
  }
\end{frame}

% ------------------------------------------------------------------
\begin{frame}{チェックポイント / Checkpoint}
  \begin{alertblock}{確認事項}
    \begin{itemize}
      \item[$\square$] 手で傾けると CF のロール・ピッチが変化する
      \item[$\square$] CF と ESKF の値が概ね一致する（Teleplot で確認）
      \item[$\square$] $\alpha$ を変えて応答の違いを観察
      \item[$\square$] Madgwick / EKF / ESKF の特徴を説明できる
    \end{itemize}
  \end{alertblock}

  \vspace{1em}

  \begin{block}{次のレッスン}
    Lesson 10: ws:: API リファレンス
  \end{block}
\end{frame}

\end{document}
